\section{Apports du portfolio pour la formation professionnelle de l'étudiante ;}
La rédaction de ce portfolio m'a appris la discipline. Devoir s'assurer d'avoir une supervision, d'anticiper les changements de stage et la perte de contact potentielle avec le superviseur. Pour être sereine lors de la réalisation du portfolio, il ne fallait pas se laisser dépasser par le stage et être alerte du cas clinique pertinent. De plus, assurer l'accès au dossier médical du patient pendant le temps nécessaire demandait de l'efficacité et une bonne gestion de son temps. Ces qualités décrites sont importantes pour un médecin et je pense qu'il s'agit d'un bon entraînement pour apprendre à compartimenter ses taches et ne pas se laisser dépasser face à différentes obligations par exemple pour la bonne gestion d'un étage.

Au niveau méthodologique, au fil du portfolio, la recherche dans la littérature a été plus simple pour moi. Au début, j'éprouvais beaucoup de difficultés d'une part à trouver des articles pertinents et maîtriser la recherche avancée. Ensuite, la lecture de l'article et déceler les potentielles erreurs ou incohérences méthodologiques des études n'était pas un exercice totalement maîtrisé pour ma part. Ce portfolio a donc été une occasion de plus pour m'entraîner et mieux saisir les subtilités de la littérature médicale.

\section{Principales questions soulevées par l’étudiante durant le portfolio en matière de compétences à acquérir ou à compléter ;}
Est-ce que la prise en charge choisit est la plus bénéfique pour le patient. En stage, confronter à l'avis de nos paires avec un symbole d'autorité et d'expérience, je prenais leur savoir comme absolue et absorbais toutes connaissances possibles. Bien sûr, grâce à la théorie apprise en cours, je pouvais être critique mais en réalité j'étais relativement naïve. Faire l'exercice de rechercher dans la littérature permet de réaliser toute la nuance de la réalité, les alternatives possibles selon le cas réel qui ne se présente pas comme dans les livres.

A travers ce travail, j'ai pu m'intéresser à l'aspect communicationnel. Soulever la question des moyens mis à disposition dans les hôpitaux pour servir au mieux la relation patient-médecin. Trouver un équilibre entre l'urgence du soin d'un côté et prendre le temps de trouver un interprète de l'autre. De plus, au-delà de la question de la temporalité, il réside toujours une question financière. Tous les hôpitaux ne disposent pas d'interprète accessible. J'ai donc pu me rendre compte de la différence entre la théorie idéale et la réalité. Je peux également ajouter la remise en question de la place du non-verbal dans la relation médecin-patient. Grâce à Yé !, je pu prendre le temps de développer cet aspect. Etre davantage alerte au langage du corps et le confronter aux paroles du patient ou bien même lorsque je suis en consultation avec une personne ne partageant pas mon langage de par son origine, une pathologie, sa culture, etc. je me concentre d'autant plus sur son non-verbal.

L'incident médical m'a permis de me positionner sur la question de l'avenir d'un enfant dans un contexte de maltraitance. Trouver un chemin de réponse pour prendre en charge ces enfants. Évaluer clairement les enjeux et les options possibles avec l'impact probable que cela aura sur l'avenir du nourrisson. Le tiraillement éthique dans ces situations est compliqué à supporter et il faut s'assurer de prendre la bonne décision.

L'application concrète du cadre légal et déontologique dans la pratique hospitalière quotidienne, en particulier dans les situations complexes, s'est révélée être un défi permanent. Entre les grands principes théoriques abordés en cours et la réalité du terrain, souvent marquée par l'urgence, le stress et la détresse humaine, l'exercice n'est jamais aisé. Prendre le temps de s'arrêter pour lire et comprendre les textes qui encadrent notre pratique (tels que la loi relative aux droits du patient ou les articles du Code de déontologie médicale) est pourtant primordial pour sécuriser et orienter nos prises en charge. L'analyse de situations cliniques délicates, notamment face à des soupçons de négligence ou de mise en danger, m'a confrontée à de réels dilemmes : comment associer l'obligation légale de protection et d'intervention avec la volonté de préserver la relation thérapeutique et la confiance avec le patient et sa famille, dans ce cas ses parents ? Ces expériences m'ont fait réaliser qu'un médecin ne peut porter seul la responsabilité de l'interprétation de la loi. A nouveau, l'équipe pluridisciplinaire a un intérêt flagrant. Se faire entourer des personnes compétentes (équipes sociales, comités d'éthique, autorités compétentes) afin de prendre une décision éclairée et collégiale n'est pas une option. 

\section{Aspects de la formation de l’étudiant·e pour lesquels le portfolio a révélé des insuffisances, et démarches entreprises pour les pallier ;}
Comprendre l'organisation d'un dossier médical. Retrouver les informations et surtout choisir les informations pertinentes. En effet, dans notre parcours académiques, les pathologies sont vues de façon dissociées, isolées l'une de l'autre or en pratique, il est commun que plusieurs pathologies co-existent au sein d'une même personne et les interactions entre elles m'étaient inconnues. Les repères sont bouleversés et la prise en charge d'autant plus complexe.

S'ajoutant au cas polypathologique est son traitement. Bien sûr, lors des études de médecin, nous acquérons des armes, des automatiques pour tenter d'adapter au mieux le traitement de personnes polymorbides. Pourtant, en pratique, ces repères ne sont pas toujours réalisable. D'une part, vu le contexte d'hospitalisation d'urgence, il n'est pas toujours possible d'initier l'arrêt de certaines molécules et de provoquer ainsi un potentiel sevrage alors que le patient n'est pas stable. Ceci pour illustrer à nouveau la différence entre les protocoles vu à l'université et la pratique concrète.

La santé mentale du soignant est clairement sous-représenté en auditoire. Nous ne sommes pas préparés à la réalité de l'hôpital. Les premiers contact avec le décès, les échecs thérapeutiques, la souffrance des patients. Malheureusement, il est encore trop souvent dit que la santé mentale et surtout ce qui pourrait être mis en place pour la préserver tel que des supports psychologiques, cohérence cardiaque, etc. ne font pas partie des « mœurs » or la détresse est réelle. Heureusement, au cours de ces cas cliniques, j'ai été en contact avec des équipes compréhensives et qui avaient la capacité de recevoir mes doutes et émotions et qui m'ont aidé à faire la part des choses.

\section{Influences que pourra avoir le Portfolio sur les perspectives professionnelles de l’étudiante}

Je pense qu'il ne faut pas chercher le médicament qui va résoudre la pathologie ou l'examen complémentaire qui va diagnostiquer le patient exactement, mais plutôt apporter de la nuance. Écouter le patient, respecter la balance bénéfice-risque pour toujours agir dans son intérêt. Pour diagnostiquer, il faudra la plupart du temps rassembler un faisceau d'arguments afin d'identifier la pathologie la plus probable. C'est dans cette même logique que je pense que le projet thérapeutique doit être élaboré.

Au-delà de la relation de soin individuelle, ce portfolio m'a également ouverte à une vision plus globale et durable de notre métier à travers l'approche \textit{One Health}. Réaliser l'interdépendance étroite entre la santé humaine, la santé animale et l'équilibre des écosystèmes — comme le démontre la problématique de l'antibiorésistance — m'a fait comprendre qu'un médecin de demain doit intégrer cette conscience environnementale et planétaire dans sa pratique.

Enfin, la littérature apporte beaucoup de guidance, mais garder le recul nécessaire, en tant que médecin, pour reprendre tous les paramètres et éléments d'un dossier garantit une réflexion profonde. L'utilisation systématique de cette posture réflexive acquise grâce au portfolio servira de véritable garde-fou professionnel. Elle me permettra d'écarter l'effet tunnel et de sécuriser mes prises de décision face à l'urgence ou à la complexité, en évitant des choix trop précipités tout au long de ma future carrière. 

\todoH[GuillaumeVertStyle][green!25]{Renvoi vers la radiographie thoracique opérationnel mais placé temporairement juste avant l'annexe A en attendant son déplacement au bon endroit dans le texte.}{Voir radiographie du thorax en \cref{annexe:radio_patient} en \cpageref{annexe:radio_patient}.}
